\subsection{Smooth support, and the four discriminants}
\label{ssec:smoothsupport}

\Cref{thm:candidatefails} and \Cref{cor:localobstruction} remove every secant
object whose support carries two transversal branches of codimension two
through a point. What they leave is a support that is smooth, or that meets
itself only along curves, and \Cref{rem:candidatemeaning} closes with the
remark that no such object is known. This subsection says why the local
argument cannot reach that case, and then determines what the case can be.
The answer is narrow: on the fourfold the invariants of the support are
forced outright, and the discriminant is confined to a finite list.

The first point is that at a smooth point there is nothing local left to
compute.

\begin{proposition}[No local obstruction at a smooth point]\label{prop:smoothlocal}
Let $Y$ be a smooth variety with $\omega_{Y}\cong\cO_{Y}$ and let $S\subset Y$
be smooth of codimension two. Then
\[
  \cE xt^{0}(\cI_{S},\cI_{S})=\cO_{Y},\qquad
  \cE xt^{1}(\cI_{S},\cI_{S})=N_{S/Y},\qquad
  \cE xt^{q}(\cI_{S},\cI_{S})=0\quad(q\ge2),
\]
and the same holds for $\cI_{S}\otimes L$ for any line bundle $L$.
Consequently, for $F=\cI_{S}\otimes L$ both sides of \eqref{eq:heisenberg}
have vanishing germs in $\cE xt^{2}$, and \Cref{lem:edge} decides nothing.
\end{proposition}

\begin{proof}
Twisting by $L$ changes neither side, so take $F=\cI_{S}$. The sheaf
$\cI_{S}$ is torsion free of rank one and reflexive, whence
$\cE xt^{0}(\cI_{S},\cI_{S})=\cO_{Y}$. Apply $\cH om(-,\cI_{S})$ to
\begin{equation}\label{eq:idealseq}
  0\to\cI_{S}\to\cO_{Y}\to\cO_{S}\to0 .
\end{equation}
Since $\cE xt^{q}(\cO_{Y},\cI_{S})=0$ for $q\ge1$, the long exact sequence
gives
\begin{equation}\label{eq:shift}
  \cE xt^{q}(\cI_{S},\cI_{S})\cong\cE xt^{q+1}(\cO_{S},\cI_{S})
  \qquad(q\ge1).
\end{equation}
Now apply $\cH om(\cO_{S},-)$ to \eqref{eq:idealseq}. Because $S$ is smooth of
codimension two, hence a local complete intersection, the Koszul resolution
gives $\cE xt^{k}(\cO_{S},\cO_{Y})=0$ for $k\ne2$ and
$\cE xt^{2}(\cO_{S},\cO_{Y})=\omega_{S}\otimes\omega_{Y}^{-1}=\omega_{S}$,
while $\cE xt^{k}(\cO_{S},\cO_{S})=\bigwedge^{k}N_{S/Y}$ for $0\le k\le2$ and
zero otherwise. Computing both from the same Koszul complex shows that the
map
\[
  \rho\colon\cE xt^{2}(\cO_{S},\cO_{Y})=\omega_{S}
  \longrightarrow
  \cE xt^{2}(\cO_{S},\cO_{S})=\textstyle\bigwedge^{2}N_{S/Y}=\omega_{S}
\]
induced by $\cO_{Y}\to\cO_{S}$ is the identity: in local coordinates in which
$S=\{x_{1}=x_{2}=0\}$, both groups are computed as $M/(x_{1},x_{2})M$ for
$M=\cO_{Y}$ and $M=\cO_{S}$, and the map is reduction modulo $(x_{1},x_{2})$,
which is the identity on $\cO_{S}$.

Take $k=2$ in the long exact sequence of $\cH om(\cO_{S},-)$. Since $\rho$ is
surjective and $\cE xt^{3}(\cO_{S},\cO_{Y})=0$, the segment
\[
  \cE xt^{2}(\cO_{S},\cO_{Y})\xrightarrow{\ \rho\ }
  \cE xt^{2}(\cO_{S},\cO_{S})\to
  \cE xt^{3}(\cO_{S},\cI_{S})\to
  \cE xt^{3}(\cO_{S},\cO_{Y})=0
\]
gives $\cE xt^{3}(\cO_{S},\cI_{S})=0$, so $\cE xt^{2}(\cI_{S},\cI_{S})=0$ by
\eqref{eq:shift}. For $k\ge3$ both $\cE xt^{k}(\cO_{S},\cO_{Y})$ and
$\cE xt^{k}(\cO_{S},\cO_{S})$ vanish, so
$\cE xt^{k+1}(\cO_{S},\cI_{S})=0$ and $\cE xt^{k}(\cI_{S},\cI_{S})=0$. Taking
$k=1$ and using that $\rho$ is injective, the segment
\[
  0=\cE xt^{1}(\cO_{S},\cO_{Y})\to
  \cE xt^{1}(\cO_{S},\cO_{S})=N_{S/Y}\to
  \cE xt^{2}(\cO_{S},\cI_{S})\to
  \cE xt^{2}(\cO_{S},\cO_{Y})\xrightarrow{\ \rho\ }
\]
exhibits $\cE xt^{2}(\cO_{S},\cI_{S})\cong N_{S/Y}$, which is
$\cE xt^{1}(\cI_{S},\cI_{S})$ by \eqref{eq:shift}.

For the last assertion, the right side of \eqref{eq:heisenberg} has vanishing
germs by \Cref{lem:edge}, and the left side lies in
$\Ext^{2}(F,F\otimes\Omega^{2}_{Y})$ whose sheaf of germs in the
$\cE xt^{2}$ direction is $\cE xt^{2}(F,F)\otimes\Omega^{2}_{Y}=0$.
\end{proof}

\begin{remark}[Where the obstruction goes instead]\label{rem:obstructionmoves}
\Cref{prop:smoothlocal} is the precise reason the argument of
\Cref{thm:candidatefails} stops at a smooth support, and it also says what
replaces it. With $\cE xt^{2}=0$ the local to global spectral sequence for
$\Ext^{2}(\cI_{S},\cI_{S})$ has only two surviving rows, $H^{1}(S,N_{S/Y})$
and $H^{2}(Y,\cO_{Y})$, and they sit in different steps of the filtration: the
first is the quotient that records the geometry of $S$, the second the
subgroup of scalars. A product $A_{\theta}A_{\theta'}$ of translation classes
therefore has a well defined image in $H^{1}(S,N_{S/Y})$, while the scalar on
the right of \eqref{eq:heisenberg} has image zero there. So the criterion
splits into two conditions of different kinds: a vanishing in
$H^{1}(S,N_{S/Y})$, which is a statement about the normal bundle and not about
any point, and an identity in $H^{2}(Y,\cO_{Y})$. At a double point the first
of these was already inconsistent locally. At a smooth support it is a global
question, and we do not decide it. What we can decide is which $S$ are even
available.
\end{remark}

The invariants of a smooth support are not free. On an abelian variety the
Todd class is trivial, so Grothendieck--Riemann--Roch turns the Chern
character of the object into the complete numerical package of the surface,
and three classical inequalities then close in on it.

\begin{theorem}[The invariants of a smooth support]\label{thm:smoothinvariants}
Let $(X,\Theta)$ be a principally polarised abelian fourfold, let $d>0$, let
$b\in\ZZ_{>0}$, and let $\cF=\cI_{S}(b\Theta)$ be a coherent sheaf with
\[
  \ch(\cF)=u_{\Theta}+b\,v_{\Theta}
\]
and $S\subset X$ smooth. Write $N=\tfrac12\Nm_{K/\QQ}(b+\sqrt{-d})
=\tfrac12(b^{2}+d)$, $\iota\colon S\hookrightarrow X$ and
$\lambda=\iota^{*}\Theta$. Then $S$ is a surface and
\begin{equation}\label{eq:smoothinvariants}
\begin{aligned}
  [S]&=N\Theta^{2}, &\qquad \iota_{*}K_{S}&=\tfrac{4bN}{3}\,\Theta^{3},
  &\qquad \chi(\cO_{S})&=4N(2b^{2}-N),\\
  K_{S}^{2}&=12N(4b^{2}-N), &\qquad e(S)&=12N(4b^{2}-3N),
  &\qquad K_{S}\cdot\lambda&=32bN,
\end{aligned}
\end{equation}
with $\lambda^{2}=24N$ and $[S]^{2}=24N^{2}$. Moreover $K_{S}$ is globally
generated, hence nef; $S$ is minimal of general type; and
\begin{equation}\label{eq:threebounds}
  \tfrac{4}{9}b^{2}\ \le\ N\ \le\ b^{2},
  \qquad\text{equivalently}\qquad
  -\tfrac{1}{9}b^{2}\ \le\ d\ \le\ b^{2} .
\end{equation}
The lower bound is the Hodge index theorem and is implied by $d>0$; the upper
bound is the Bogomolov--Miyaoka--Yau inequality. In terms of $b$ and $d$ the
five invariants read
\begin{equation}\label{eq:bdform}
  \chi(\cO_{S})=(b^{2}+d)(3b^{2}-d),\quad
  K_{S}^{2}=3(b^{2}+d)(7b^{2}-d),\quad
  e(S)=3(b^{2}+d)(5b^{2}-3d),
\end{equation}
and the defect in the Bogomolov--Miyaoka--Yau inequality is
\begin{equation}\label{eq:bmydefect}
  3e(S)-K_{S}^{2}=24\,(b^{4}-d^{2}),
\end{equation}
so that the inequality is exactly $d\le b^{2}$ and nothing else.
Finally $K_{S}$ is never numerically proportional to $\lambda$.
\end{theorem}

\begin{proof}
Since $\ch_{0}(\cF)=1$ and $c_{1}(\cF)=b\Theta$, the sheaf
$\cF(-b\Theta)=\cI_{S}$ has $\ch_{0}=1$, $\ch_{1}=0$, so $S$ has codimension
at least two, and
\[
  \ch(\cO_{S})=1-\ch(\cI_{S})=1-\ch(\cF)\,e^{-b\Theta}.
\]
Expanding the right side in $\QQ[\Theta]/(\Theta^{5})$ and collecting gives
\[
  \ch_{2}(\cO_{S})=\tfrac{b^{2}+d}{2}\Theta^{2}=N\Theta^{2},\quad
  \ch_{3}(\cO_{S})=-\tfrac{2bN}{3}\Theta^{3},\quad
  \ch_{4}(\cO_{S})=\tfrac{3b^{4}+2db^{2}-d^{2}}{24}\Theta^{4},
\]
which is the computation of item (XXX) of \Cref{sec:verification}. The first
identity says $S$ has codimension exactly two, so it is a surface, with class
$N\Theta^{2}$. Because $T_{X}$ is trivial, $\td(T_{X})=1$, and
Grothendieck--Riemann--Roch for $\iota$ reads
$\ch(\iota_{*}\cO_{S})=\iota_{*}\td(T_{S})$. Comparing degrees three and four,
\[
  \ch_{3}(\cO_{S})=\iota_{*}\tfrac{c_{1}(T_{S})}{2}=-\tfrac12\iota_{*}K_{S},
  \qquad
  \ch_{4}(\cO_{S})=\iota_{*}\tfrac{c_{1}^{2}+c_{2}}{12}
  =\chi(\cO_{S})\,[\mathrm{pt}],
\]
the second by Noether's formula. With $\Theta^{4}=24\,[\mathrm{pt}]$ these give
$\iota_{*}K_{S}=\tfrac{4bN}{3}\Theta^{3}$ and, after substituting
$d=2N-b^{2}$, $\chi(\cO_{S})=3b^{4}+2db^{2}-d^{2}=4N(2b^{2}-N)$.
Then $K_{S}\cdot\lambda=\int_{X}\iota_{*}K_{S}\cdot\Theta=\tfrac{4bN}{3}\cdot24=32bN$
and $\lambda^{2}=\int_{X}[S]\cdot\Theta^{2}=24N$.

For the last two invariants, the normal bundle sequence and the triviality of
$T_{X}$ give $c(T_{S})c(N_{S/X})=1$, so $c_{1}(N_{S/X})=K_{S}$ and
$c_{2}(N_{S/X})=K_{S}^{2}-e(S)$; the self-intersection formula
$[S]^{2}=\int_{S}c_{2}(N_{S/X})$ then reads
\[
  24N^{2}=K_{S}^{2}-e(S),
\]
while Noether's formula reads $12\chi(\cO_{S})=K_{S}^{2}+e(S)$. Solving the
two gives the stated $K_{S}^{2}$ and $e(S)$.

The conormal sequence presents $\Omega^{1}_{S}$ as a quotient of
$\Omega^{1}_{X}|_{S}=\cO_{S}^{\oplus4}$, so $\Omega^{1}_{S}$ is globally
generated, and therefore so is $K_{S}=\det\Omega^{1}_{S}$, a quotient of
$\bigwedge^{2}\cO_{S}^{\oplus4}$. A globally generated line bundle is nef. A
globally generated rank two bundle on a surface has $c_{2}\ge0$, the class
$c_{2}$ being represented by the vanishing scheme of a general section, which
is empty or finite; hence
\[
  e(S)=c_{2}(\Omega^{1}_{S})=12N(4b^{2}-3N)\ \ge\ 0,
\]
so $N\le\tfrac43b^{2}<4b^{2}$ and therefore $K_{S}^{2}=12N(4b^{2}-N)>0$. A nef
divisor of positive self-intersection is big, so $\kappa(S)=2$ and, $K_{S}$
being nef, $S$ is minimal of general type. The Bogomolov--Miyaoka--Yau
inequality $c_{1}^{2}\le3c_{2}$ then applies and reads
$K_{S}^{2}\le3e(S)$. Substituting $N=\tfrac12(b^{2}+d)$ in the two
expressions gives \eqref{eq:bdform}, hence
$3e(S)-K_{S}^{2}=24(b^{4}-d^{2})$, and the inequality is $d\le b^{2}$, that
is $N\le b^{2}$.

For the lower bound, $\lambda$ is ample on $S$, so the Hodge index theorem
gives $(K_{S}\cdot\lambda)^{2}\ge K_{S}^{2}\,\lambda^{2}$, that is
$(32bN)^{2}\ge12N(4b^{2}-N)\cdot24N$, which simplifies to $9N\ge4b^{2}$. Since
$N=\tfrac12(b^{2}+d)$ with $d>0$, this holds automatically. Equality in the
Hodge index theorem occurs exactly when $K_{S}$ is numerically proportional to
$\lambda$, hence exactly when $9N=4b^{2}$; as $9N=\tfrac92(b^{2}+d)>4b^{2}$ for
$d>0$, that never happens.
\end{proof}

\begin{corollary}[Four discriminants on the fourfold]\label{cor:fourdiscriminants}
In \Cref{setup:markman}, where $b=3$ by \Cref{prop:markman8}(iii), a secant
object $\cI_{S}(3\Theta)$ with smooth support exists only if $d\le9$; since
$d$ is squarefree this leaves $d\in\{1,2,3,5,6,7\}$, and if in addition
$N=(9+d)/2$ is an integer, which is what a support assembled from translates
of $\Theta$ requires, only
\[
  d\in\{1,\,3,\,5,\,7\},\qquad N\in\{5,\,6,\,7,\,8\}.
\]
For these the invariants are forced to the values in
\Cref{tab:smoothsupport}. In particular, for every imaginary quadratic field
of discriminant larger than $9$ there is no secant object of rank one with
smooth support on an abelian fourfold at all.
\end{corollary}

\begin{proof}
Immediate from \eqref{eq:threebounds} at $b=3$, together with the arithmetic
of \eqref{eq:smoothinvariants}; the table is item (XXX) of
\Cref{sec:verification}. The value $d=9$ is the equality case of
Bogomolov--Miyaoka--Yau, where $K_{S}^{2}=3e(S)=2916$ and $S$ would have to be
a quotient of the ball; it is excluded here because $9$ is not squarefree.
\end{proof}

\begin{table}[htbp]
\centering
\small
\renewcommand{\arraystretch}{1.12}
\begin{tabular}{@{}ccrrrrr@{}}
\toprule
$d$ & $N$ & $\chi(\cO_{S})$ & $K_{S}^{2}$ & $e(S)$ &
$K_{S}\cdot\lambda$ & $\lambda^{2}$ \\
\midrule
$1$ & $5$ & $260$ & $1860$ & $1260$ & $480$ & $120$ \\
$3$ & $6$ & $288$ & $2160$ & $1296$ & $576$ & $144$ \\
$5$ & $7$ & $308$ & $2436$ & $1260$ & $672$ & $168$ \\
$7$ & $8$ & $320$ & $2688$ & $1152$ & $768$ & $192$ \\
\bottomrule
\end{tabular}
\caption{The only invariants a smooth support can have on an abelian
fourfold, by \Cref{thm:smoothinvariants} and
\Cref{cor:fourdiscriminants}, at $b=3$ and with $N$ an integer. Every entry is
forced by $\ch(\cF)$ alone; none is a choice. The slope
$K_{S}^{2}/\chi(\cO_{S})$ runs from $93/13$ at $d=1$ to $42/5$ at $d=7$, and
reaches the Bogomolov--Miyaoka--Yau ceiling of $9$ only at the excluded value
$d=9$.}
\label{tab:smoothsupport}
\end{table}

\begin{figure}[htbp]
\centering
  \includegraphics[width=\textwidth]{fig_bmywindow}
\caption{\Cref{thm:smoothinvariants} over the quarter plane of $(b,d)$. The
height of the surface is the defect $3e(S)-K_{S}^{2}=24(b^{4}-d^{2})$ in the
Bogomolov--Miyaoka--Yau inequality for the surface the theorem attaches to the
pair, compressed by a signed cube root so that four orders of magnitude fit
one frame; the compression is odd and monotone, so the zero locus and the sign
are the true ones. The surface is drawn indigo where the defect is positive and
clay where it is negative, and it crosses zero along $d=b^{2}$, which is
therefore the whole content of the inequality. The weaker wall $d=\tfrac53
b^{2}$, where $e(S)$ itself changes sign, is drawn dashed and lies above it.
On the line $b=3$, which is the value \Cref{prop:markman8}(iii) forces on the
fourfold, the four squarefree discriminants $d=1,3,5,7$ with $N$ an integer sit
above the wall and are marked with their drop lines; $d=9$ sits exactly on it,
the equality case, and $d=11,13,15$ sit below and are excluded.}
\label{fig:bmywindow}
\end{figure}

\begin{remark}[What is left of the smooth case]\label{rem:smoothremains}
Three things have changed and one has not.

The local argument is out of reach here, and \Cref{prop:smoothlocal} says so
for a structural reason rather than as a failure of effort: the sheaf
$\cE xt^{2}$ of a smooth ideal vanishes, so both sides of
\eqref{eq:heisenberg} are locally zero and the germ comparison that decides
\Cref{thm:candidatefails} has nothing to compare.

The surface is pinned. Its class, its canonical class, its Euler
characteristic, its Euler number and its two intersection numbers against the
polarisation are all determined by $d$ alone, with no freedom, and they are
compatible with existence only for $d\le9$. That is the content of
\Cref{cor:fourdiscriminants}, and it is a genuine restriction: the secant
construction with smooth support is not available for all but finitely many
imaginary quadratic fields, whatever geometry one is willing to invent.

The residual question is now a single existence statement. For
$d\in\{1,3,5,7\}$, does a principally polarised abelian fourfold carry a smooth
surface $S$ with the invariants of \Cref{tab:smoothsupport}, and if so does
the image of $A_{\theta}A_{\theta'}$ in $H^{1}(S,N_{S/X})$ vanish for all
$\theta,\theta'\in H^{0}(T_{X})$? The first half is a question about surfaces
in abelian fourfolds with prescribed Chern numbers, of the kind that is
answered by construction or by a classification argument; the second is a
question about one cohomology group of one normal bundle. Neither is settled
here.

What has not changed is the standing of the conjecture. Even a positive answer
to both would produce Weil classes on abelian eightfolds for at most four
discriminants, and \Cref{cor:exactform} would still be open for every other
discriminant, for every $n\ge5$, and for every field of degree larger than
two.
\end{remark}

\subsection{The obstruction is a condition on depth}
\label{ssec:depth}

\Cref{prop:smoothlocal} says that at a smooth point of the support there is
nothing local to compare. That is a special case of something much more
general, and identifying the general case tells us exactly which supports
\Cref{thm:candidatefails} and \Cref{cor:localobstruction} remove and which
they cannot touch. The dividing line is not smoothness, and it is not being a
local complete intersection. It is depth.

\begin{theorem}[Cohen--Macaulay supports carry no local obstruction]%
\label{thm:cmvanishing}
Let $Y$ be a smooth variety with $\omega_{Y}\cong\cO_{Y}$, let $Z\subset Y$ be
a closed subscheme of codimension two, and let $p\in Z$. If the local ring
$\cO_{Z,p}$ is Cohen--Macaulay, then
\[
  \cE xt^{q}(\cI_{Z},\cI_{Z})_{p}=0\qquad\text{for every }q\ge2 ,
\]
and the same holds for $\cI_{Z}\otimes L$ for every line bundle $L$.
Consequently both sides of \eqref{eq:heisenberg} have vanishing germ at $p$,
and \Cref{lem:edge} decides nothing there.
\end{theorem}

\begin{proof}
Twisting by $L$ changes no $\cE xt$ sheaf, so take $F=\cI_{Z}$. Write
$A=\cO_{Y,p}$, a regular local ring of dimension $m$, and $I=\cI_{Z,p}$. By
hypothesis $A/I$ is Cohen--Macaulay of dimension $m-2$, so
$\operatorname{depth}A/I=m-2$, and the Auslander--Buchsbaum formula gives
\[
  \operatorname{pd}_{A}(A/I)=m-\operatorname{depth}_{A}(A/I)=2 .
\]
From the exact sequence $0\to I\to A\to A/I\to0$ with $A$ free it follows that
$\operatorname{pd}_{A}I=1$. A module of projective dimension one has a free
resolution of length one, so $\Ext^{q}_{A}(I,M)=0$ for every $A$-module $M$
and every $q\ge2$; taking $M=I$ gives the vanishing. The last assertion is
\Cref{lem:edge}: the right side of \eqref{eq:heisenberg} is
$z\cdot\id_{F}$ with $z\in H^{2}(\cO_{Y})$ and has vanishing germs, and the
left side lies in a sheaf whose $\cE xt^{2}$ factor is zero.
\end{proof}

\begin{remark}[The hypothesis is sharp, and it is the one that fails]%
\label{rem:cmsharp}
Two observations make \Cref{thm:cmvanishing} the right statement.

The first is that Cohen--Macaulayness is strictly weaker than the conditions
one might reach for instead. Item (XXXI) of \Cref{sec:verification} computes
$\cE xt^{2}(\cI_{Z},\cI_{Z})$ for ten codimension two germs in $\CC^{4}$. It
vanishes for the smooth germ $(x_{1},x_{2})$; for two planes meeting along a
line, $(x_{1},x_{2}x_{3})$, and for three planes meeting along a line,
$(x_{1},x_{2}x_{3}x_{4})$; for the double plane $(x_{1}^{2},x_{2})$ and the
quadric cone $(x_{1},x_{2}^{2}-x_{3}x_{4})$; and also for the fat plane
$(x_{1},x_{2})^{2}$ and for the cone over the twisted cubic, the ideal of
$2\times2$ minors of $\left(\begin{smallmatrix}x_{1}&x_{2}&x_{3}\\
x_{2}&x_{3}&x_{4}\end{smallmatrix}\right)$, neither of which is a local
complete intersection, both needing three generators in codimension two. So
reducedness, irreducibility and the complete intersection property are all
irrelevant; the projective dimension is what matters, and by
Auslander--Buchsbaum that is depth.

The second is that the germs where it does not vanish are exactly the ones
the construction of \Cref{setup:markman} produces. The same computation finds
$\cE xt^{2}\ne0$, with nonzero products of translation classes, for two planes
meeting at a point, for three planes meeting pairwise at a point, and for a
plane with an embedded point; the projective dimensions there are $3$, $3$ and
$4$. A scheme whose local ring is Cohen--Macaulay is connected in codimension
one at that point, by Hartshorne's connectedness theorem, so a support with
two components through $p$ meeting only at $p$ can never be Cohen--Macaulay
there. That is precisely the configuration \Cref{lem:localproducts} analyses.
\end{remark}

\begin{corollary}[What the local argument removes, in final form]%
\label{cor:cmdichotomy}
Let $(X,t)$ be a polarised abelian variety of dimension $n\ge4$ and $F$ a
perfect complex with $\ch(F)=au_{t}+bv_{t}$, $(a,b)\ne(0,0)$, whose germ at
each point is the germ of $\cI_{Z}\otimes L$ for a codimension two subscheme
$Z$. Then exactly one of the following holds.
\begin{enumerate}[label=\textup{(\roman*)},leftmargin=2.6em]
\item $Z$ fails to be Cohen--Macaulay at some point $p$ of one of the types of
  \Cref{rem:cmsharp}, and then $F$ does not satisfy the weakened criterion,
  nor does $\Phi(F\boxtimes F_{2})$ for any secant $F_{2}$, by
  \Cref{thm:candidatefails} and \Cref{cor:localobstruction}.
\item $Z$ is Cohen--Macaulay, and then there is no local obstruction at any
  point whatever. The criterion is then a single global condition: the image
  of $A_{\theta}A_{\theta'}$ in $H^{1}\bigl(Z,\cE xt^{1}(\cI_{Z},\cI_{Z})\bigr)$
  must vanish for all $\theta,\theta'\in H^{0}(T_{X})$, the identity in
  $H^{2}(\cO_{X})$ being all that the scalar on the right of
  \eqref{eq:heisenberg} can meet.
\end{enumerate}
In case (ii), $\cI_{Z}$ has projective dimension one at every point, so it is
locally the ideal of maximal minors of a Hilbert--Burch matrix, and when $Z$
is in addition smooth $\cE xt^{1}(\cI_{Z},\cI_{Z})=N_{Z/X}$ and the invariants
of $Z$ are those of \Cref{thm:smoothinvariants}.
\end{corollary}

\begin{proof}
The two cases are exclusive and exhaustive by definition. In case (i) the
local ring is not Cohen--Macaulay and the cited results apply, the germs being
those of \Cref{lem:localproducts} and \Cref{cor:localobstruction}. In case
(ii) \Cref{thm:cmvanishing} gives $\cE xt^{q}=0$ for $q\ge2$ at every point, so
the local to global spectral sequence for $\Ext^{2}(\cI_{Z},\cI_{Z})$ has only
the two rows $H^{1}(\cE xt^{1})$ and $H^{2}(\cO_{X})$, and these lie in
different steps of the filtration, the second being the subgroup of scalars.
The right side of \eqref{eq:heisenberg} lies in that subgroup, so it has zero
image in the quotient $H^{1}(\cE xt^{1})$, and the criterion forces the image
of the left side there to vanish. The projective dimension statement is the
proof of \Cref{thm:cmvanishing}, and Hilbert--Burch is the structure theorem
for modules of projective dimension one over a regular local ring; the last
clause is \Cref{prop:smoothlocal} and \Cref{thm:smoothinvariants}.
\end{proof}

\begin{remark}[Why the normalisation cannot repair the candidate]%
\label{rem:normalisationcm}
\Cref{cor:cmdichotomy} explains the tension in \Cref{rem:normalisedpoints} in
one word. The partial normalisation performed in \Cref{setup:markman}
separates the two branches at an isolated double point, and separating them is
exactly what would make the local ring Cohen--Macaulay there. So the
construction is already trying to reach case (ii). What
\Cref{rem:normalisedpoints} shows is that it cannot arrive: normalising every
isolated point moves $\chi(\cO_{\widetilde Z})$ to $50N$, while the secant
plane demands $(d+9)(27-d)=72N-4N^{2}$, and the two agree only at $N=11/2$.
Repairing the depth destroys the Chern character, and keeping the Chern
character keeps a point where the depth fails. That is the obstruction in its
final form, and it is why a construction assembled from intersections of
translates of a polarisation cannot be made to work: by
\Cref{cor:localobstruction} such a chain always has a point of the bad type,
and by \Cref{rem:normalisedpoints} it cannot be separated without leaving the
plane.
\end{remark}

\subsection{Cohen--Macaulay supports, and the two routes that remain}
\label{ssec:hilbertburch}

\Cref{cor:cmdichotomy} leaves one class of supports and one condition on them.
The class is the Cohen--Macaulay subschemes of codimension two, and for those
the proof of \Cref{thm:cmvanishing} gives more than the vanishing it was used
for: the ideal sheaf has projective dimension one at every point. That is the
Hilbert--Burch shape, and it can be searched.

\begin{setupenv}[A resolution of the support]\label{setup:burch}
Let $(X,\Theta)$ be a principally polarised abelian fourfold with
$\NS(X)_{\QQ}=\QQ\Theta$, let $Z\subset X$ be Cohen--Macaulay of codimension
two, and suppose $\cI_{Z}$ admits a global resolution
\begin{equation}\label{eq:burchres}
  0\longrightarrow E_{1}\stackrel{\varphi}{\longrightarrow}E_{0}
  \longrightarrow\cI_{Z}\longrightarrow0,
  \qquad \rk E_{0}=r+1,\quad\rk E_{1}=r,
\end{equation}
with $E_{0}=\bigoplus_{i=1}^{r+1}\cO(e_{i}\Theta)$ and
$E_{1}=\bigoplus_{j=1}^{r}\cO(f_{j}\Theta)$ sums of line bundles. Put
$p_{i}=e_{i}+b$ and $q_{j}=f_{j}+b$.
\end{setupenv}

\begin{lemma}[The resolution in moments]\label{lem:burchmoments}
In \Cref{setup:burch}, $\ch(\cI_{Z}(b\Theta))=u_{\Theta}+b\,v_{\Theta}$ holds
if and only if
\begin{equation}\label{eq:burchsystem}
  \sum_{i=1}^{r+1}p_{i}^{\,k}-\sum_{j=1}^{r}q_{j}^{\,k}=c_{k},
  \qquad k=0,1,2,3,4,
  \qquad (c_{k})=(1,\;b,\;-d,\;-db,\;d^{2}).
\end{equation}
The equation at $k=0$ is the rank count and is automatic. Moreover
$\varphi$ can be injective only if $\max_{i}e_{i}\ge\max_{j}f_{j}$, and
$\cO_{X}$ embeds in $E_{0}^{\vee}$ only if $\min_{i}e_{i}\le0$.
\end{lemma}

\begin{proof}
From \eqref{eq:burchres}, $\ch(\cI_{Z})=\ch(E_{0})-\ch(E_{1})$, so
$\ch(\cI_{Z}(b\Theta))=\sum_{i}e^{p_{i}\Theta}-\sum_{j}e^{q_{j}\Theta}$.
Comparing coefficients of $\Theta^{k}/k!$ against $u_{\Theta}+bv_{\Theta}$,
whose coefficients are the $c_{k}$ displayed, gives
\eqref{eq:burchsystem}. For the last two statements, on an abelian variety
$\Hom(\cO(f\Theta),\cO(e\Theta))=H^{0}(\cO((e-f)\Theta))$ vanishes unless
$e\ge f$; if $\varphi$ is injective then each summand of $E_{1}$ must map
nonzero to $E_{0}$, which needs $\max_{i}e_{i}\ge f_{j}$ for every $j$.
Dualising \eqref{eq:burchres} and using $\cI_{Z}^{\vee}=\cO_{X}$ gives an
injection $\cO_{X}\hookrightarrow E_{0}^{\vee}$, which needs
$-e_{i}\ge0$ for some $i$.
\end{proof}

The case $r=1$ is the complete intersections, and it closes outright, for
every $b$ and every $d$, by an argument that uses no search.

\begin{theorem}[No complete intersection supports a secant object]%
\label{thm:noci}
Let $(X,\Theta)$ be a principally polarised abelian fourfold with
$\NS(X)_{\QQ}=\QQ\Theta$, let $d>0$ and $b\ge1$. No complete intersection
$Z=D_{1}\cap D_{2}$ of two divisors satisfies
$\ch(\cI_{Z}(b\Theta))=u_{\Theta}+b\,v_{\Theta}$.
\end{theorem}

\begin{proof}
Such a $Z$ is \Cref{setup:burch} with $r=1$, the resolution being the Koszul
complex of the two equations. Put $w(x)=x^{2}+d$, which is strictly positive
on $\RR$ because $d>0$, and read \eqref{eq:burchsystem} in the combinations
$c_{k+2}+d\,c_{k}$ for $k=0,1,2$. Those three numbers are
\[
  c_{2}+dc_{0}=-d+d=0,\qquad
  c_{3}+dc_{1}=-db+db=0,\qquad
  c_{4}+dc_{2}=d^{2}-d^{2}=0,
\]
so that
\[
  \sum_{i}w(p_{i})\,p_{i}^{\,k}=\sum_{j}w(q_{j})\,q_{j}^{\,k},
  \qquad k=0,1,2 .
\]
The two sides are the moments of order $0$, $1$ and $2$ of the positive
measures $\nu_{P}=\sum_{i}w(p_{i})\delta_{p_{i}}$ and
$\nu_{Q}=\sum_{j}w(q_{j})\delta_{q_{j}}$, so those measures have the same
mass, the same mean and the same variance. At $r=1$ the measure $\nu_{Q}$ has
one atom, hence variance zero, so $\nu_{P}$ has variance zero as well, and a
positive measure with two atoms has variance zero only if the atoms coincide.
So $p_{1}=p_{2}$, and then equality of mass and mean forces $q_{1}=p_{1}$. The
resolution \eqref{eq:burchres} then cancels and $\cI_{Z}\cong\cO(e_{1}\Theta)$
is a line bundle, which the ideal sheaf of a subscheme of codimension two is
not.

Equivalently and by the same computation, the two divisor degrees $a_{1}$ and
$a_{2}$ would have to satisfy $a_{1}+a_{2}=\tfrac43b$ and
$(b-a_{1})(b-a_{2})=\tfrac12\bigl(\tfrac13b^{2}+d\bigr)$, and the quadratic
with that sum and product has discriminant $-\tfrac29b^{2}-2d<0$: the degrees
are not even real.
\end{proof}

For $r\ge2$ the moment conditions no longer decide by themselves, and
\eqref{eq:burchsystem} becomes a finite search over the twists. Item (XXXII)
of \Cref{sec:verification} carries it out exactly, by meet in the middle on
the four moments, and the outcome is \Cref{tab:burch}.

\begin{theorem}[The split resolutions, and where they can live]%
\label{thm:burchsearch}
In \Cref{setup:burch} with $b=3$, over every squarefree $d<24$ and every
choice of twists in the box recorded in \Cref{tab:burch}:
\begin{enumerate}[label=\textup{(\roman*)},leftmargin=2.6em]
\item at $r=2$ and $r=3$ the system \eqref{eq:burchsystem} has no solution at
  all, and the same holds for every $b\le6$ over seven fields;
\item at $r=4,5,6$ it has solutions, but only for $d=15$ and $d=23$;
\item every $d$ that survives at any rank exceeds $9$.
\end{enumerate}
\end{theorem}

\begin{proof}
The search is item (XXXII) of \Cref{sec:verification}, carried out in exact
integer arithmetic; the two conditions of \Cref{lem:burchmoments} on
$\max_{i}e_{i}$ and $\min_{i}e_{i}$ are imposed, and without the first the
system acquires solutions that are not resolutions, the smallest being
$p=(-8,-5,0,1,2)$, $q=(-7,-6,-4,4)$ at $r=4$ and $d=23$, where the summand of
$E_{1}$ with the largest twist admits no nonzero map to $E_{0}$.
\end{proof}

\begin{table}[htbp]
\centering
\small
\renewcommand{\arraystretch}{1.12}
\begin{tabular}{@{}clll@{}}
\toprule
$r$ & twists searched & $d$ surviving & first candidate, as twists of $\Theta$ \\
\midrule
$1$ & every $b$, every $d$ & none & \Cref{thm:noci} \\
$2$ & $[-18,18]$ & none & \\
$3$ & $[-14,14]$ & none & \\
$4$ & $[-11,11]$ & $23$ & $E_{0}:(-11,-6,0,2,7)$,\ \ $E_{1}:(-10,-9,5,6)$ \\
$5$ & $[-9,9]$ & $23$ & $E_{0}:(-8,-4,-3,-2,1,6)$,\ \ $E_{1}:(-7,-6,-6,4,5)$ \\
$6$ & $[-8,8]$ & $15,23$ & $E_{0}:(-8,-4,-4,-1,-1,1,4)$,\ \ $E_{1}:(-7,-6,-6,0,3,3)$ \\
\bottomrule
\end{tabular}
\caption{The split Hilbert--Burch resolutions of \Cref{setup:burch} at $b=3$,
over every squarefree $d<24$. Nothing survives below rank four, and what
survives above it lives only at $d=15$ and $d=23$. Both exceed the bound $d\le9$
of \Cref{cor:fourdiscriminants}, so none of these supports can be smooth.}
\label{tab:burch}
\end{table}


The search of \Cref{thm:burchsearch} records where a split resolution can
have the right Chern character. It is superseded by a statement that needs no
search and no bound on the twists: a split resolution never runs the
criterion, at any rank and any discriminant. The reason is a necessary
condition that is worth isolating, because it applies to every object and can
be tested on any candidate before any global computation is attempted.

\begin{lemma}[Morphisms to a line bundle are killed by $H^{2}(\cO)$]
\label{lem:annihilation}
Let $(X,t)$ be a polarised abelian variety of dimension $n\ge3$, let $F$ be
a perfect complex with $\ch(F)=au_{t}+bv_{t}$, $(a,b)\ne(0,0)$, which
satisfies the weakened criterion of \Cref{setup:criterion}, and let $L$ be a
line bundle with $c_{1}(L)\in\QQ\,t$. Then for every $m\in\ZZ$, every
morphism $p\colon F\to L[m]$, every morphism $q\colon L[m]\to F$, and every
$z\in H^{2}(\cO_{X})$,
\[
  z\cdot p=0\ \ \text{in }\Hom(F,L[m+2]),\qquad
  z\cdot q=0\ \ \text{in }\Hom(L[m],F[2]),
\]
where $z\cdot(\,\cdot\,)$ is the $H^{*}(\cO_{X})$-module structure on
morphisms in $D^{b}(X)$. In particular, if $F=I_{Z}\otimes L'$ for a closed
subscheme $Z$ and a line bundle $L'$ with $c_{1}(L')\in\QQ t$, then the
restriction $H^{n-2}(\cO_{X})\to H^{n-2}(\cO_{Z})$ is injective.
\end{lemma}

\begin{proof}
For a constant bivector $\pi$ the class
$x_{\pi}=(\tfrac d2\,\pi\lrcorner\,t^{2},0,\pi)$ preserves $\ch(F)$ by
\Cref{thm:factor}(ii), so $\operatorname{ev}_{F}(x_{\pi})=0$ by hypothesis.
The Hochschild action is by natural transformations $\id\to\id[2]$ of
$D^{b}(X)$, so for $p\colon F\to L[m]$,
$\operatorname{ev}_{L[m]}(x_{\pi})\circ p=p\circ\operatorname{ev}_{F}(x_{\pi})=0$,
and $\operatorname{ev}_{L[m]}=\operatorname{ev}_{L}[m]$. By \eqref{eq:evsplit},
$\operatorname{ev}_{L}(x_{\pi})=c_{\pi}\cdot\id_{L}$ with
$c_{\pi}=\tfrac12\,\pi\lrcorner\,(\lambda^{2}+d\,t^{2})\in H^{0,2}(X)$,
$\lambda=c_{1}(L)=e\,t$, that is $c_{\pi}=\tfrac{e^{2}+d}{2}\,\pi\lrcorner\,t^{2}$,
and $\pi\mapsto\pi\lrcorner\,t^{2}$ is an isomorphism from
$H^{0}(\bigwedge^{2}T_{X})$ onto $H^{0,2}(X)$, as in the proof of
\Cref{thm:factor}(iii). Hence $c_{\pi}$ runs over all of $H^{2}(\cO_{X})$ as
$\pi$ varies, and $(c_{\pi}\cdot\id_{L})\circ p=c_{\pi}\cdot p$. The
statement for $q$ is the same with the square
$\operatorname{ev}_{F}(x_{\pi})\circ q=q\circ\operatorname{ev}_{L[m]}(x_{\pi})$.
For the last assertion apply the first to the inclusion
$p\colon I_{Z}\otimes L'\to L'$: the sequence
$0\to I_{Z}\to\cO_{X}\to\cO_{Z}\to0$ gives an exact sequence
\[
  \Ext^{2}(\cO_{Z},\cO_{X})\longrightarrow H^{2}(\cO_{X})
  \longrightarrow\Ext^{2}(I_{Z},\cO_{X}),
\]
in which the second map is $z\mapsto z\cdot p$; so $z\cdot p=0$ for all $z$
means that the first map is surjective, and by Serre duality with $\omega_{X}$
trivial it is the dual of the restriction
$H^{n-2}(\cO_{X})\to H^{n-2}(\cO_{Z})$.
\end{proof}

\begin{theorem}[No two-term complex of powers of the polarisation runs the
criterion]\label{thm:nosplit}
Let $(X,t)$ be a polarised abelian variety of dimension $n\ge3$ and $M$ an
ample line bundle with $c_{1}(M)\in\QQ t$. Let $A$ and $B$ be direct sums of
tensor powers $M^{\otimes e}$, $e\in\ZZ$, let $\varphi\colon A\to B$ be any
morphism, and let $F$ be the cone of $\varphi$, or any shift of it. If $F\ne0$
and $\ch(F)=au_{t}+bv_{t}$ with $(a,b)\ne(0,0)$, then $F$ does not satisfy
the weakened criterion. In particular no object with a resolution
$0\to E_{1}\to E_{0}\to F\to0$ by sums of powers of $M$, whatever the rank,
does.
\end{theorem}

\begin{proof}
A pair of summands $M^{\otimes e}$ of $A$ and $M^{\otimes e}$ of $B$ joined
by a nonzero entry of $\varphi$, which is then a nonzero scalar, can be
cancelled without changing the cone up to isomorphism; do this until no such
pair remains. If $A=0$ or $B=0$ after cancelling, $F$ is a shifted sum of
line bundles, $\operatorname{ev}_{F}(x_{\pi})$ is diagonal with the nonzero
entries $c_{\pi}\cdot\id$ of the proof of \Cref{lem:annihilation}, and the
criterion fails as in \Cref{cor:noline}; so suppose both are nonzero and
consider the triangle $A\xrightarrow{\varphi}B\xrightarrow{\;q\;}F\to A[1]$.
Let $M^{\otimes b_{j}}$ be a summand of $B$, $\iota_{j}$ its inclusion into
$B$, and $q_{j}=q\circ\iota_{j}\colon M^{\otimes b_{j}}\to F$. Applying
$\Hom(M^{\otimes b_{j}},-)$ to the triangle gives the exact sequence
\[
  \Hom(M^{\otimes b_{j}},A[2])\xrightarrow{\;\varphi_{*}\;}
  \Hom(M^{\otimes b_{j}},B[2])\xrightarrow{\;q_{*}\;}
  \Hom(M^{\otimes b_{j}},F[2]),
\]
and $z\cdot q_{j}=q_{*}(z\cdot\iota_{j})$ for $z\in H^{2}(\cO_{X})$. Now
$\Hom(M^{\otimes b_{j}},B[2])=\bigoplus_{k}H^{2}(M^{\otimes(b_{k}-b_{j})})$,
and $H^{2}(M^{\otimes m})=0$ for $m\ne0$: for $m>0$ by Kodaira vanishing,
$\omega_{X}$ being trivial, and for $m<0$ because a negative line bundle on
an abelian variety has cohomology only in degree $n\ge3$. So $z\cdot\iota_{j}$
is the element with component $z$ in the summand $k=j$ and zero elsewhere. It
lies in the image of $\varphi_{*}$ only if there is
$\psi\in\Hom(M^{\otimes b_{j}},A[2])=\bigoplus_{i}H^{2}(M^{\otimes(a_{i}-b_{j})})$
with $(\varphi\circ\psi)_{j}=\sum_{i}\varphi_{ji}\psi_{i}=z$; but $\psi_{i}\ne0$
forces $a_{i}=b_{j}$, and for such $i$ the entry $\varphi_{ji}$ is a scalar,
zero after the cancellation. So $z\cdot q_{j}\ne0$ for every $z\ne0$, which
contradicts \Cref{lem:annihilation}. The last assertion is the case of a
resolution, where $F\cong\coker\varphi$ is the cone of $\varphi\colon E_{1}\to E_{0}$.
\end{proof}

The powers of the polarisation are not the only bundles on which the
compensated Poisson classes act by scalars. So does every bundle whose Atiyah
class is a multiple of the identity, and on an abelian variety these are the
semi-homogeneous bundles of Mukai \cite{Muk78}; the argument of
\Cref{thm:nosplit} goes through for them, with the cohomology of the fibre
product of two isogenies in place of the cohomology of a power of $\Theta$.
We state the result in a form that is self-contained.

\begin{definition}[Isogeny blocks]\label{def:block}
Let $(X,\Theta)$ be a polarised abelian variety. An \emph{isogeny block} is a
vector bundle $E=\pi_{*}L$, where $\pi\colon Y\to X$ is an isogeny and $L$ is
a line bundle on $Y$ with $c_{1}(L)=\nu\,\pi^{*}\Theta$ in $H^{2}(Y,\QQ)$,
$\nu\in\QQ$; its \emph{slope} is $\nu$, so that $c_{1}(E)=\nu\deg(\pi)\,\Theta$.
The block is \emph{simple} if $\tau_{g}^{*}L\not\cong L$ for every nonzero
$g\in\ker\pi$. The powers $\cO(e\Theta)$ are the blocks with $\pi=\id$.
\end{definition}

\begin{lemma}[Three properties of a block]\label{lem:blocks}
Let $E=\pi_{*}L$ be an isogeny block of slope $\nu$ on an abelian $n$-fold.
\begin{enumerate}[label=\textup{(\roman*)},nosep]
\item $\mathrm{At}(E)=\nu\Theta\cdot\id_{E}$; hence for a constant bivector
  $\pi$ the compensated class $x_{\pi}$ of \Cref{thm:factor}(ii) acts by
  $\operatorname{ev}_{E}(x_{\pi})=\tfrac{\nu^{2}+d}{2}\,(\pi\lrcorner\,\Theta^{2})\cdot\id_{E}$,
  and \Cref{lem:annihilation} holds with $E$ in place of the line bundle $L$.
\item For a second block $E'=\pi'_{*}L'$ of slope $\nu'$,
  $\Ext^{q}(E',E)=H^{q}(Z,M)$ with $Z=Y'\times_{X}Y$, a disjoint union of
  abelian varieties isogenous to $X$, and $M$ a line bundle whose class on
  each component is $(\nu-\nu')$ times the pullback of $\Theta$. If
  $\nu\ne\nu'$ and $n\ge3$ then $\Ext^{2}(E',E)=0$. If $E$ is simple then
  $\Ext^{2}(E,E)=H^{2}(\cO_{X})\cdot\id_{E}$.
\item A simple block is $\mu$-stable with respect to $\Theta$; hence a
  nonzero morphism between simple blocks of the same slope is an
  isomorphism.
\end{enumerate}
\end{lemma}

\begin{proof}
(i) Since $\pi$ is \'etale, $\pi^{*}\Omega^{1}_{X}=\Omega^{1}_{Y}$ and $\pi_{*}$
is exact, so the first jet sequence of $\pi_{*}L$ is the direct image of
that of $L$, and $\mathrm{At}(\pi_{*}L)$ is the image of $\mathrm{At}(L)$
under $\pi_{*}\colon\Ext^{1}_{Y}(L,L\otimes\pi^{*}\Omega^{1}_{X})\to
\Ext^{1}_{X}(\pi_{*}L,\pi_{*}L\otimes\Omega^{1}_{X})$. As $\mathrm{At}(L)=c_{1}(L)\cdot\id_{L}
=\nu\,\pi^{*}\Theta\cdot\id_{L}$, the projection formula gives
$\nu\Theta\cdot\id_{\pi_{*}L}$. Then \eqref{eq:evgeneral} with a scalar Atiyah
class gives the value of $\operatorname{ev}_{E}(x_{\pi})$, and the proof of
\Cref{lem:annihilation} used nothing else about $L$.

(ii) Adjunction for the finite \'etale $\pi'$, flat base change along the
cartesian square with vertices $Z$, $Y'$, $Y$, $X$, and the projection formula
give $\Ext^{q}_{X}(\pi'_{*}L',\pi_{*}L)=H^{q}(Y',L'^{-1}\otimes\pi'^{*}\pi_{*}L)
=H^{q}(Z,p'^{*}L'^{-1}\otimes p^{*}L)$, with $p,p'$ the two projections; the
class of $M=p'^{*}L'^{-1}\otimes p^{*}L$ is $(\nu-\nu')$ times the pullback of
$\Theta$ on every component, which is an abelian variety finite \'etale over
$Y'$. A line bundle on an abelian $n$-fold whose class is a positive
multiple of an ample class has cohomology only in degree $0$, and one whose
class is a negative multiple only in degree $n$, by Kodaira vanishing and
Serre duality with trivial canonical bundle; so $H^{2}(Z,M)=0$ when
$\nu\ne\nu'$ and $n\ge3$. When $E'=E$, the components of $Y\times_{X}Y$ are
the graphs of the translations by $g\in\ker\pi$, on which $M$ restricts to
$L^{-1}\otimes\tau_{g}^{*}L\in\Pic^{0}(Y)$; a nontrivial element of $\Pic^{0}$
has no cohomology, so $\Ext^{q}(E,E)=\bigoplus_{g}H^{q}(\cO_{Y})$ over the
$g$ with $\tau_{g}^{*}L\cong L$, which for a simple block is $g=0$ alone, and
$H^{q}(\cO_{Y})=H^{q}(\cO_{X})$ through $\pi$.

(iii) The pullback $\pi^{*}E=\bigoplus_{g\in\ker\pi}\tau_{g}^{*}L$ is a direct
sum of line bundles of one slope, hence $\mu$-semistable, and semistability
descends along the finite \'etale $\pi$; so $E$ is semistable. If $F\subset E$
is a saturated subsheaf of the same slope, then $\pi^{*}F\subset\pi^{*}E$ is
a saturated subsheaf of the same slope of a direct sum of pairwise
non-isomorphic line bundles, the $\tau_{g}^{*}L$ being distinct because the
block is simple, hence a partial sum $\bigoplus_{g\in S}\tau_{g}^{*}L$; it is
invariant under the translations by $\ker\pi$, which permute the summands
transitively, so $S$ is empty or everything and $F=0$ or $F=E$. A nonzero
morphism between stable bundles of the same slope is an isomorphism.
\end{proof}

\begin{theorem}[No two-term complex of simple blocks runs the criterion]
\label{thm:noblocks}
Let $(X,\Theta)$ be a polarised abelian variety of dimension $n\ge3$, let $A$
and $B$ be direct sums of simple isogeny blocks, let $\varphi\colon A\to B$
be any morphism and let $F$ be the cone of $\varphi$ or a shift of it. If
$F\ne0$ and $\ch(F)=au_{\Theta}+bv_{\Theta}$ with $(a,b)\ne(0,0)$, then $F$
does not satisfy the weakened criterion. In particular no object with a
resolution $0\to E_{1}\to E_{0}\to F\to0$ by sums of simple isogeny blocks
does, whatever the rank.
\end{theorem}

\begin{proof}
Cancel, by Gaussian elimination in the two-term complex, every pair of
isomorphic blocks of $A$ and $B$ joined by an entry of $\varphi$ that is an
isomorphism; this changes neither the cone nor the block structure of the
remaining terms. By \Cref{lem:blocks}(iii), every remaining entry between
blocks of the same slope is zero. The rest is the proof of
\Cref{thm:nosplit} with \Cref{lem:blocks} in place of the vanishing of
$H^{2}(M^{\otimes m})$: for a block $E'$ of $B$ with inclusion $\iota$ and
$q'=q\circ\iota\colon E'\to F$, the class $z\cdot\iota\in\Hom(E',B[2])$ has
component $z\cdot\id_{E'}\ne0$ in $\Ext^{2}(E',E')=H^{2}(\cO_{X})\cdot\id_{E'}$
by (ii); it lies in the image of $\varphi_{*}$ only if
$\sum_{i}\varphi_{ji}\psi_{i}=z\cdot\id_{E'}$ for some
$\psi_{i}\in\Ext^{2}(E',A_{i})$, and $\psi_{i}\ne0$ forces $A_{i}$ to have
the slope of $E'$ by (ii), for which $\varphi_{ji}=0$. So $z\cdot q'\ne0$,
against \Cref{lem:annihilation} in the form of (i). If $A$ or $B$ is empty
after cancelling, $F$ is a shifted sum of blocks, on each of which
$\operatorname{ev}(x_{\pi})$ is a nonzero scalar by (i).
\end{proof}

\begin{remark}[What the blocks are]\label{rem:blocks}
By Mukai \cite{Muk78} the simple isogeny blocks are exactly the simple
semi-homogeneous bundles, the bundles $E$ with $\tau_{x}^{*}E\cong E\otimes P$
for every $x$ and some $P\in\Pic^{0}(X)$, and every semi-homogeneous bundle
on a variety with $\NS(X)_{\QQ}=\QQ\Theta$ is an iterated extension of such
blocks and their twists by $\Pic^{0}$; we do not use this. What
\Cref{thm:noblocks} says is that the bundles in a Hilbert--Burch resolution
of a candidate cannot all have scalar Atiyah class. The obstruction is the
same one every time: the compensated Poisson classes act on a block by the
scalar $\tfrac{\nu^{2}+d}{2}\,\pi\lrcorner\,\Theta^{2}$, which would vanish only
for $\nu^{2}=-d$, and cones of maps between blocks inherit the scalar through
the summand inclusions because the $\Ext^{2}$ between blocks of different
slopes vanishes. A candidate must therefore be built from a bundle whose
Atiyah class is not a scalar, that is, one which is not projectively flat.
\end{remark}

\begin{corollary}[The split shape is closed, and what the two routes must be]
\label{cor:disjointroutes}
Let $(X,\Theta)$ be a principally polarised abelian fourfold.
\begin{enumerate}[label=\textup{(\roman*)},nosep]
\item No secant object of rank one whose support has a resolution
  \eqref{eq:burchres} by sums of powers of $\Theta$ satisfies the weakened
  criterion, for any $r$, any twists and any $d$; in particular every
  candidate of \Cref{tab:burch} is excluded, whether or not its matrix
  degenerates in codimension two.
\item If a secant object $\cI_{Z}(b\Theta)$ satisfies the criterion, then
  $Z$ is Cohen--Macaulay, $\cI_{Z}$ has projective dimension one, its
  Hilbert--Burch resolution does not split into powers of $\Theta$, nor into
  simple semi-homogeneous bundles, and the restriction
  $H^{2}(\cO_{X})\to H^{2}(\cO_{Z})$ is injective. If $Z$ is smooth then
  moreover $d\in\{1,3,5,7\}$.
\item In the range of \Cref{thm:burchsearch} the two shapes do not meet:
  a smooth support forces $d\le9$ and a split resolution forces $d\ge15$,
  so even before (i) no support could be both.
\end{enumerate}
\end{corollary}

\begin{proof}
(i) is \Cref{thm:nosplit} with $M=\cO_{X}(\Theta)$, the resolution
\eqref{eq:burchres} presenting $\cI_{Z}(b\Theta)$ as the cone of
$\varphi\otimes\cO(b\Theta)$. (ii): Cohen--Macaulay and projective
dimension one are \Cref{cor:cmdichotomy}(ii); non-splitting is (i) and
\Cref{thm:noblocks};
injectivity is the last assertion of \Cref{lem:annihilation}; the four
discriminants are \Cref{cor:fourdiscriminants}. (iii) is
\Cref{cor:fourdiscriminants} against \Cref{thm:burchsearch}(iii), the sets
$\{1,3,5,7\}$ and $\{15,23\}$ being disjoint.
\end{proof}

\begin{remark}[The route map]\label{rem:routemap}
Collecting what is now known about the first of the six statements of
\Cref{ssec:closuremap}, the position is this. The local part of the criterion
is finished: it is a condition on depth, nothing else
(\Cref{thm:cmvanishing}, \Cref{cor:cmdichotomy}). The support must therefore
be Cohen--Macaulay, its ideal has projective dimension one, and there are
exactly two shapes left.

The first is a smooth support. Then every invariant is forced, the
discriminant is one of four (\Cref{cor:fourdiscriminants}), and the resolution
cannot split, now for every $d$ and every rank (\Cref{thm:nosplit}). What is
needed is the existence of
such a surface on a principally polarised abelian fourfold, with a
Hilbert--Burch matrix over a bundle that is not a sum of line bundles, and
then the vanishing of the image of $A_{\theta}A_{\theta'}$ in
$H^{1}(S,N_{S/X})$. The first half is a question about surfaces with
prescribed Chern numbers, of the kind settled by construction or by
classification; the second is a question about one cohomology group.

The second is a singular Cohen--Macaulay support. Then the invariants are
still forced by \eqref{eq:smoothinvariants} with $c_{1}(N_{Z/X})$ in place of
$K_{Z}$, and by \Cref{cor:disjointroutes}(i) the Hilbert--Burch resolution
must not split into powers of $\Theta$ or into simple semi-homogeneous
bundles (\Cref{thm:noblocks}): the explicit candidates of \Cref{tab:burch}
are excluded, and what is needed is a resolution over a bundle that is not
projectively flat, together with the vanishing in
$H^{1}(Z,\cE xt^{1}(\cI_{Z},\cI_{Z}))$. In both shapes the restriction
$H^{2}(\cO_{X})\to H^{2}(\cO_{Z})$ must be injective
(\Cref{lem:annihilation}), which for a smooth $S$ says that the six classes
of $\bigwedge^{2}H^{1}(\cO_{X})$ remain independent in $H^{2}(\cO_{S})$.

Both are finite questions in the sense that they are questions about one
variety and one sheaf, not about a family or a limit. Neither is settled here,
and settling both would give $\mathbf{W}(K,4,\delta)$ for the discriminants
reached and nothing beyond them.
\end{remark}
